\scanpage{7}
des $(\varphi_i)$ (``structure Kummérienne'')
compte tenu des opérations \inserted{(16) et (18)} de $\mathcal G_{X'}$
sur $I$ et $\mathbf T(\overline{\mathbf Q})$. On peut donc
écrire
\[
(19)\qquad\mathcal G_{X'}\longrightarrow\operatorname{Autext}(\pi_1(X'),\mathbf T(\overline{\mathbf Q}),(\varphi_i)).
\]
\deleted{Supposons $\unsure{X}$}

Soit $\widehat{\widetilde X'}$ le normalisé de $\deleted{\unsure{\widetilde{\phantom X}}}X$ dans ``l'extension
\unsure{étale|totale}'' définie par $\widetilde X'$, donc on a un
morphisme profini
\[
\widehat{\widetilde X'}\longrightarrow X
\]
Soit $\widetilde X'(I)$ l'image inverse de l'ens des
pts \deleted{génériques} \inserted{minimaux} de $D=\sum D_i$. On a
donc
\[
(20)\qquad\widetilde X'(I)\mathbin{-}I
\]
et une opération de $\deleted{\unsure{\mathcal G_{X'}}}\mathcal E_{X'}$ sur
$\widetilde X'(I)$ (par transport de structure)
compatible avec les opérations sur $I$
et l'application (20). Pour $i\in I$ soit
\[
(21)\qquad\widetilde X'(i)\subset\widetilde X'(I)
\]
sa fibre dans $\widetilde X'(I)$, \inserted{(c'est un espace hom.
sous $\pi_1(X')$. On trouve alors}
\mspage{6}
